\chapter{Returns, Variance, and Why Volatility Matters}
\label{ch:returns}

\begin{application}[Why This Chapter]
Volatility is the central quantity in quantitative finance.
Options pricing, risk management, portfolio construction, and trade execution all depend on accurate volatility estimates.
Every chapter in this guide builds on the concepts introduced here.
All five project directions (Chapter~\ref{ch:applications}) start from the returns and variance foundations covered in this chapter.
\end{application}

Financial markets produce a stream of prices.
To do anything quantitative with those prices, you first need to convert them into \emph{returns}, measure how much those returns vary, and understand \emph{why} that variation matters.
This chapter covers that ground.

% ═══════════════════════════════════════════════════
\section{What Are Returns?}
\label{sec:returns-definition}
% ═══════════════════════════════════════════════════

A stock price on its own tells you almost nothing about performance.
Knowing that Apple closed at \$187 today is useless without knowing where it closed yesterday.
\emph{Returns} solve this: they express price changes as fractions (or percentages), making different assets and time periods comparable.

There are two standard ways to compute a return.

\begin{prereq}[Natural Logarithm Properties]
The natural logarithm $\ln(\cdot)$ is the inverse of the exponential function: $\ln(e^x) = x$.
Three properties matter here:
\begin{itemize}[nosep]
  \item $\ln(A/B) = \ln A - \ln B$ (quotients become differences)
  \item $\ln(AB) = \ln A + \ln B$ (products become sums)
  \item For small $x$: $\ln(1+x) \approx x$ (so log returns $\approx$ simple returns when returns are small)
\end{itemize}
\end{prereq}

\subsection{Simple Returns}
\label{sec:simple-returns}

The simple (arithmetic) return over one period measures the fractional price change:

\begin{equation}
  R_t = \frac{P_t - P_{t-1}}{P_{t-1}}
  \label{eq:simple-return}
\end{equation}
\begin{itemize}[nosep]
  \item $R_t$: simple return from period $t-1$ to $t$
  \item $P_t$: price at the end of period $t$
  \item $P_{t-1}$: price at the end of the previous period
\end{itemize}

A simple return of $0.05$ means the price rose by 5\%.

\subsection{Log Returns}
\label{sec:log-returns}

The log (continuously compounded) return uses the natural logarithm of the price ratio:

\begin{equation}
  r_t = \ln\!\left(\frac{P_t}{P_{t-1}}\right) = \ln P_t - \ln P_{t-1}
  \label{eq:log-return}
\end{equation}
\begin{itemize}[nosep]
  \item $r_t$: log return from period $t-1$ to $t$
  \item $\ln(\cdot)$: natural logarithm
\end{itemize}

\begin{keyidea}[Why Log Returns?]
Log returns have two properties that make them the default in volatility research:
\begin{enumerate}[nosep]
  \item \textbf{Time additivity.}
  The multi-period log return is the sum of single-period log returns:
  $r_{1:T} = r_1 + r_2 + \cdots + r_T$.
  Simple returns compound multiplicatively, which is harder to work with.
  \item \textbf{Approximate symmetry.}
  A log return of $+0.05$ followed by $-0.05$ returns you close to your starting price.
  Simple returns do not: a $+5\%$ gain followed by a $-5\%$ loss leaves you at $\$99.75$, not $\$100$.
\end{enumerate}
Throughout this guide, $r_t$ always means the log return unless stated otherwise.
\end{keyidea}

\begin{workedexample}{Computing Returns}
A stock closes at \$100.00 on Monday, \$105.00 on Tuesday, and \$102.00 on Wednesday.

\textbf{Tuesday (price rises from \$100 to \$105):}
\begin{align*}
  R_{\text{Tue}} &= \frac{105 - 100}{100} = 0.0500 \quad (+5.00\%) \\[4pt]
  r_{\text{Tue}} &= \ln\!\left(\frac{105}{100}\right) = \ln(1.05) = 0.0488 \quad (+4.88\%)
\end{align*}

\textbf{Wednesday (price falls from \$105 to \$102):}
\begin{align*}
  R_{\text{Wed}} &= \frac{102 - 105}{105} = -0.0286 \quad (-2.86\%) \\[4pt]
  r_{\text{Wed}} &= \ln\!\left(\frac{102}{105}\right) = \ln(0.9714) = -0.0290 \quad (-2.90\%)
\end{align*}

\textbf{Two-day log return (time additivity in action):}
\[
  r_{\text{Mon}\to\text{Wed}} = r_{\text{Tue}} + r_{\text{Wed}} = 0.0488 + (-0.0290) = 0.0198
\]
Check directly: $\ln(102/100) = \ln(1.02) = 0.0198$. \checkmark

The simple returns are close to the log returns because the magnitudes are small (under 5\%).
For daily equity returns, which are typically well under 1\%, the two are nearly identical.
\end{workedexample}

% ═══════════════════════════════════════════════════
\section{Variance and Standard Deviation}
\label{sec:variance}
% ═══════════════════════════════════════════════════

With returns in hand, the next question is: how spread out are they?
Variance answers this by measuring the average squared deviation from the mean return.

\begin{definition}[Sample Variance and Standard Deviation]
Given $T$ log returns $r_1, r_2, \ldots, r_T$ with sample mean $\bar{r} = \frac{1}{T}\sum_{t=1}^T r_t$, the sample variance is:

\begin{equation}
  \hat{\sigma}^2 = \frac{1}{T-1}\sum_{t=1}^{T}\bigl(r_t - \bar{r}\bigr)^2
  \label{eq:sample-variance}
\end{equation}
\begin{itemize}[nosep]
  \item $\hat{\sigma}^2$: sample variance (the hat denotes an estimate)
  \item $T$: number of observations
  \item $r_t$: log return at time $t$
  \item $\bar{r}$: sample mean of returns
  \item $T-1$: Bessel's correction (dividing by $T-1$ instead of $T$ gives an unbiased estimate)
\end{itemize}

The sample standard deviation is $\hat{\sigma} = \sqrt{\hat{\sigma}^2}$.
\end{definition}

\begin{prereq}[Why $T-1$?]
Dividing by $T-1$ instead of $T$ corrects for the fact that you estimated $\bar{r}$ from the same data.
Using $\bar{r}$ instead of the true mean $\mu$ systematically shrinks the deviations, so dividing by the smaller number $T-1$ compensates.
This is called Bessel's correction.
For large $T$ the difference is negligible.
\end{prereq}

\subsection{Annualizing Volatility}
\label{sec:annualizing}

Raw daily standard deviations are tiny numbers (around 0.01 for a typical equity index), so practitioners report volatility on an annual scale.
If daily returns are independent, the variance over $n$ days is $n$ times the one-day variance, so the standard deviation scales by $\sqrt{n}$.

\begin{equation}
  \hat{\sigma}_{\text{annual}} = \hat{\sigma}_{\text{daily}} \times \sqrt{252}
  \label{eq:annualize}
\end{equation}
\begin{itemize}[nosep]
  \item $252$: approximate number of trading days per year in U.S.\ equity markets
  \item $\sqrt{252} \approx 15.87$
\end{itemize}

\begin{warning}[The Square-Root-of-Time Rule]
The $\sqrt{252}$ scaling assumes returns are independent and identically distributed.
When volatility clusters (Section~\ref{sec:stylized-facts}), this assumption fails, and multi-day volatility may be higher or lower than the square-root rule predicts.
Chapters~\ref{ch:garch} and \ref{ch:har} develop models that handle this.
\end{warning}

\begin{workedexample}{From Daily Returns to Annualized Volatility}
Five consecutive daily log returns for the S\&P~500 (in decimal):
\[
  r_1 = 0.0032,\quad r_2 = -0.0051,\quad r_3 = 0.0018,\quad r_4 = 0.0074,\quad r_5 = -0.0028
\]

\textbf{Step 1: Mean.}
\[
  \bar{r} = \frac{0.0032 + (-0.0051) + 0.0018 + 0.0074 + (-0.0028)}{5} = \frac{0.0045}{5} = 0.0009
\]

\textbf{Step 2: Squared deviations.}
\begin{align*}
  (r_1 - \bar{r})^2 &= (0.0032 - 0.0009)^2 = (0.0023)^2 = 0.00000529 \\
  (r_2 - \bar{r})^2 &= (-0.0051 - 0.0009)^2 = (-0.0060)^2 = 0.00003600 \\
  (r_3 - \bar{r})^2 &= (0.0018 - 0.0009)^2 = (0.0009)^2 = 0.00000081 \\
  (r_4 - \bar{r})^2 &= (0.0074 - 0.0009)^2 = (0.0065)^2 = 0.00004225 \\
  (r_5 - \bar{r})^2 &= (-0.0028 - 0.0009)^2 = (-0.0037)^2 = 0.00001369
\end{align*}

\textbf{Step 3: Variance.}
\[
  \hat{\sigma}^2 = \frac{0.00000529 + 0.00003600 + 0.00000081 + 0.00004225 + 0.00001369}{4}
               = \frac{0.00009804}{4}
               = 0.00002451
\]

\textbf{Step 4: Daily and annualized volatility.}
\[
  \hat{\sigma}_{\text{daily}} = \sqrt{0.00002451} = 0.00495 \quad (0.50\%\text{ per day})
\]
\[
  \hat{\sigma}_{\text{annual}} = 0.00495 \times \sqrt{252} = 0.00495 \times 15.87 = 0.0786 \quad (7.9\%)
\]

For context, the S\&P~500's long-run annualized volatility is roughly 15--20\%.
Our five-day sample yields a lower figure because five observations cannot capture the full range of market conditions.
\end{workedexample}

% ═══════════════════════════════════════════════════
\section{Why Volatility Matters}
\label{sec:why-vol-matters}
% ═══════════════════════════════════════════════════

You now know how to compute returns and measure their spread.
The next question: why does that spread matter so much?
Variance and standard deviation sit at the center of four major problems in finance.

\begin{prereq}[Finance Concepts Preview]
The following four applications are covered briefly here to motivate the rest of the guide.
You do not need to understand them fully yet; each connects to a later chapter.
\end{prereq}

\paragraph{Options pricing.}
An option gives the holder the right (not obligation) to buy or sell an asset at a fixed price by a future date.
The value of this right depends critically on how much the underlying price might move, which is volatility.
The Black-Scholes model \citep{Black1973} takes volatility as an \emph{input} and produces an option price as output.
If your volatility estimate is wrong, your price is wrong.
Chapter~\ref{ch:volsurface} explores this connection.

\paragraph{Risk management.}
Value-at-Risk ($\VaR$) estimates the worst-case loss over a holding period at a given confidence level.
In the simplest (Gaussian) version, $\VaR$ is proportional to portfolio volatility:
$\VaR_\alpha = -\mu + z_\alpha \,\sigma$,
where $z_\alpha$ is the normal quantile.
When volatility doubles, the risk bound roughly doubles.

\paragraph{Portfolio construction.}
Mean-variance optimization \citep{Markowitz1952} and risk-parity strategies both require a covariance matrix as input.
Volatility is the diagonal of that matrix.
Better volatility forecasts produce better-diversified portfolios.

\paragraph{Trade execution.}
When a desk needs to buy a large position, it splits the order over time to reduce market impact.
The participation rate (fraction of daily volume per time slice) depends on intraday volatility: higher volatility means wider price swings, which changes optimal execution speed.

\begin{keyidea}[What Makes Volatility Special]
Volatility is one of the few quantities in finance that is (a)~directly observable from intraday data (Chapter~\ref{ch:rv}), (b)~economically central to pricing, hedging, and risk control, and (c)~forecastable with meaningful accuracy (Chapters~\ref{ch:garch}--\ref{ch:hybrid}).
That combination makes it the right place to apply ML carefully.
\end{keyidea}

% ═══════════════════════════════════════════════════
\section{Stylized Facts of Financial Returns}
\label{sec:stylized-facts}
% ═══════════════════════════════════════════════════

Before building models, you need to know what patterns actually appear in return data.
\citet{Cont2001} catalogued a set of ``stylized facts'': statistical regularities observed across equities, foreign exchange, and commodities, across decades and geographies.
Four of these facts matter most for volatility modeling.

\subsection{Fact 1: Returns Are Approximately Uncorrelated}
\label{sec:fact-uncorrelated}

Daily return autocorrelations are statistically indistinguishable from zero for lags beyond a few minutes \citep{Cont2001}.
In plain terms: knowing today's return tells you almost nothing about tomorrow's return.
This is consistent with market efficiency; if returns were predictable, traders would exploit the pattern until it disappeared.

\begin{intuition}[No Free Lunch in Means]
If positive returns reliably followed positive returns, everyone would buy after an up day, driving the price up immediately and eliminating the pattern.
Competition among traders keeps return autocorrelations near zero.
Volatility, by contrast, is \emph{not} a quantity you can directly trade on the same way, so its autocorrelation can persist.
\end{intuition}

\subsection{Fact 2: Volatility Clusters}
\label{sec:fact-clustering}

Although returns themselves are uncorrelated, \emph{squared} returns and \emph{absolute} returns show strong, slowly decaying positive autocorrelation \citep{Cont2001, Mandelbrot1963}.
Large moves tend to follow large moves, and calm periods tend to follow calm periods.

\citet{Engle1982} formalized this observation with the ARCH model: the variance of next period's return depends on recent squared returns.
Chapter~\ref{ch:garch} develops this family of models in detail.

\begin{keyresult}[\citet{Cont2001}, Stylized Fact: Volatility Clustering]
The autocorrelation of absolute returns remains significantly positive over lags of several weeks and decays slowly (approximately as a power law with exponent $\beta \in [0.2, 0.4]$), a pattern that is remarkably stable across asset classes and time periods.
\end{keyresult}

Figure~\ref{fig:vol-clustering-ts} shows volatility clustering visually: returns alternate between calm and turbulent stretches.

\begin{figure}[htbp]
\centering
\begin{tikzpicture}
\begin{axis}[
  width=\textwidth,
  height=6cm,
  xlabel={Trading day},
  ylabel={Daily log return (\%)},
  xmin=0, xmax=500,
  ymin=-5, ymax=5,
  grid=major,
  grid style={gray!20},
  tick label style={font=\small},
  label style={font=\small},
  title style={font=\small\bfseries},
  title={Simulated Daily Returns Exhibiting Volatility Clustering},
]
% Calm regime (days 1-100): low vol ~0.5%
\addplot[defblue, mark=none, thick, opacity=0.8] coordinates {
(1,0.12)(2,-0.31)(3,0.45)(4,-0.18)(5,0.27)(6,0.08)(7,-0.42)(8,0.33)(9,-0.15)(10,0.22)
(11,-0.38)(12,0.19)(13,0.05)(14,-0.28)(15,0.41)(16,-0.11)(17,0.35)(18,-0.22)(19,0.14)(20,-0.06)
(21,0.29)(22,-0.33)(23,0.17)(24,-0.41)(25,0.23)(26,0.09)(27,-0.36)(28,0.28)(29,-0.13)(30,0.44)
(31,-0.21)(32,0.16)(33,-0.38)(34,0.31)(35,-0.09)(36,0.25)(37,-0.19)(38,0.37)(39,-0.27)(40,0.11)
(41,0.03)(42,-0.34)(43,0.26)(44,-0.15)(45,0.39)(46,-0.22)(47,0.18)(48,-0.41)(49,0.32)(50,-0.08)
(51,0.21)(52,-0.29)(53,0.13)(54,0.36)(55,-0.24)(56,0.07)(57,-0.33)(58,0.28)(59,-0.17)(60,0.42)
(61,-0.11)(62,0.25)(63,-0.38)(64,0.19)(65,-0.06)(66,0.31)(67,-0.23)(68,0.14)(69,-0.35)(70,0.27)
(71,0.04)(72,-0.29)(73,0.21)(74,-0.16)(75,0.38)(76,-0.12)(77,0.33)(78,-0.25)(79,0.09)(80,-0.41)
(81,0.22)(82,-0.07)(83,0.35)(84,-0.28)(85,0.15)(86,-0.33)(87,0.26)(88,-0.18)(89,0.40)(90,-0.14)
(91,0.08)(92,-0.31)(93,0.24)(94,-0.19)(95,0.36)(96,-0.10)(97,0.29)(98,-0.37)(99,0.16)(100,-0.23)
};
% Turbulent regime (days 101-220): high vol ~2.5%
\addplot[warnred, mark=none, thick, opacity=0.8] coordinates {
(101,-1.82)(102,2.45)(103,-2.91)(104,1.67)(105,-3.12)(106,2.78)(107,-1.95)(108,3.41)(109,-2.53)(110,1.89)
(111,-3.28)(112,2.14)(113,-1.73)(114,2.96)(115,-2.41)(116,1.52)(117,-3.05)(118,2.63)(119,-1.88)(120,3.17)
(121,-2.72)(122,1.94)(123,-2.36)(124,3.08)(125,-1.61)(126,2.49)(127,-2.87)(128,1.75)(129,-3.33)(130,2.21)
(131,-1.56)(132,2.84)(133,-2.19)(134,1.43)(135,-2.95)(136,2.57)(137,-1.81)(138,3.22)(139,-2.64)(140,1.98)
(141,-2.11)(142,2.73)(143,-3.45)(144,1.67)(145,-2.38)(146,2.91)(147,-1.79)(148,3.14)(149,-2.52)(150,1.86)
(151,-2.97)(152,2.33)(153,-1.64)(154,2.78)(155,-3.21)(156,1.91)(157,-2.46)(158,2.65)(159,-1.73)(160,3.08)
(161,-2.89)(162,2.17)(163,-1.58)(164,2.94)(165,-2.31)(166,1.76)(167,-3.15)(168,2.48)(169,-1.92)(170,2.71)
(171,-2.54)(172,1.83)(173,-3.07)(174,2.39)(175,-1.66)(176,2.86)(177,-2.23)(178,1.57)(179,-2.98)(180,2.61)
(181,-1.84)(182,3.19)(183,-2.67)(184,1.95)(185,-2.42)(186,2.74)(187,-1.71)(188,3.11)(189,-2.55)(190,1.88)
(191,-3.24)(192,2.16)(193,-1.59)(194,2.93)(195,-2.37)(196,1.72)(197,-3.09)(198,2.52)(199,-1.87)(200,2.68)
(201,-2.18)(202,1.54)(203,-2.81)(204,2.36)(205,-1.69)(206,2.97)(207,-2.49)(208,1.83)(209,-3.16)(210,2.24)
(211,-1.62)(212,2.79)(213,-2.34)(214,1.91)(215,-2.73)(216,2.47)(217,-1.78)(218,3.05)(219,-2.58)(220,1.86)
};
% Calm regime (days 221-350)
\addplot[defblue, mark=none, thick, opacity=0.8] coordinates {
(221,-0.28)(222,0.34)(223,-0.17)(224,0.41)(225,-0.23)(226,0.09)(227,-0.36)(228,0.28)(229,-0.14)(230,0.32)
(231,-0.21)(232,0.15)(233,-0.39)(234,0.26)(235,-0.08)(236,0.43)(237,-0.31)(238,0.18)(239,-0.24)(240,0.37)
(241,-0.12)(242,0.29)(243,-0.35)(244,0.22)(245,-0.06)(246,0.38)(247,-0.27)(248,0.13)(249,-0.41)(250,0.25)
(251,-0.19)(252,0.33)(253,-0.16)(254,0.44)(255,-0.29)(256,0.11)(257,-0.37)(258,0.24)(259,-0.08)(260,0.36)
(261,-0.22)(262,0.17)(263,-0.34)(264,0.28)(265,-0.13)(266,0.42)(267,-0.26)(268,0.09)(269,-0.38)(270,0.31)
(271,-0.15)(272,0.27)(273,-0.33)(274,0.19)(275,-0.07)(276,0.41)(277,-0.25)(278,0.14)(279,-0.36)(280,0.23)
(281,-0.18)(282,0.35)(283,-0.29)(284,0.12)(285,-0.43)(286,0.26)(287,-0.11)(288,0.38)(289,-0.22)(290,0.16)
(291,-0.31)(292,0.24)(293,-0.09)(294,0.37)(295,-0.28)(296,0.13)(297,-0.35)(298,0.21)(299,-0.06)(300,0.39)
(301,-0.23)(302,0.17)(303,-0.34)(304,0.28)(305,-0.14)(306,0.41)(307,-0.26)(308,0.11)(309,-0.38)(310,0.22)
(311,-0.19)(312,0.33)(313,-0.16)(314,0.44)(315,-0.29)(316,0.08)(317,-0.36)(318,0.25)(319,-0.12)(320,0.37)
(321,-0.21)(322,0.15)(323,-0.32)(324,0.27)(325,-0.09)(326,0.43)(327,-0.24)(328,0.18)(329,-0.35)(330,0.23)
(331,-0.17)(332,0.31)(333,-0.14)(334,0.39)(335,-0.27)(336,0.12)(337,-0.36)(338,0.24)(339,-0.08)(340,0.42)
(341,-0.23)(342,0.16)(343,-0.33)(344,0.29)(345,-0.11)(346,0.38)(347,-0.25)(348,0.13)(349,-0.37)(350,0.22)
};
% Turbulent regime (days 351-420)
\addplot[warnred, mark=none, thick, opacity=0.8] coordinates {
(351,-1.42)(352,2.15)(353,-2.63)(354,1.87)(355,-2.28)(356,1.54)(357,-2.91)(358,2.36)(359,-1.73)(360,2.84)
(361,-2.17)(362,1.62)(363,-2.79)(364,2.43)(365,-1.88)(366,2.11)(367,-2.56)(368,1.79)(369,-3.01)(370,2.27)
(371,-1.65)(372,2.72)(373,-2.34)(374,1.91)(375,-2.48)(376,2.16)(377,-1.83)(378,2.69)(379,-2.25)(380,1.57)
(381,-2.87)(382,2.38)(383,-1.74)(384,2.93)(385,-2.41)(386,1.68)(387,-2.76)(388,2.14)(389,-1.59)(390,2.81)
(391,-2.33)(392,1.86)(393,-2.67)(394,2.42)(395,-1.78)(396,2.55)(397,-2.19)(398,1.63)(399,-2.88)(400,2.31)
(401,-1.71)(402,2.64)(403,-2.27)(404,1.94)(405,-2.53)(406,2.18)(407,-1.82)(408,2.76)(409,-2.39)(410,1.61)
(411,-2.84)(412,2.29)(413,-1.67)(414,2.91)(415,-2.46)(416,1.85)(417,-2.72)(418,2.13)(419,-1.58)(420,2.47)
};
% Return to calm (days 421-500)
\addplot[defblue, mark=none, thick, opacity=0.8] coordinates {
(421,-0.24)(422,0.31)(423,-0.18)(424,0.42)(425,-0.27)(426,0.13)(427,-0.35)(428,0.22)(429,-0.09)(430,0.38)
(431,-0.21)(432,0.16)(433,-0.33)(434,0.28)(435,-0.14)(436,0.41)(437,-0.25)(438,0.11)(439,-0.37)(440,0.23)
(441,-0.19)(442,0.34)(443,-0.15)(444,0.29)(445,-0.36)(446,0.17)(447,-0.26)(448,0.39)(449,-0.12)(450,0.32)
(451,-0.22)(452,0.14)(453,-0.38)(454,0.27)(455,-0.08)(456,0.43)(457,-0.31)(458,0.19)(459,-0.24)(460,0.36)
(461,-0.16)(462,0.28)(463,-0.34)(464,0.21)(465,-0.07)(466,0.41)(467,-0.23)(468,0.15)(469,-0.32)(470,0.26)
(471,-0.11)(472,0.38)(473,-0.25)(474,0.17)(475,-0.36)(476,0.22)(477,-0.13)(478,0.33)(479,-0.29)(480,0.08)
(481,-0.37)(482,0.24)(483,-0.18)(484,0.42)(485,-0.26)(486,0.12)(487,-0.35)(488,0.28)(489,-0.14)(490,0.31)
(491,-0.21)(492,0.16)(493,-0.39)(494,0.23)(495,-0.09)(496,0.34)(497,-0.27)(498,0.18)(499,-0.33)(500,0.25)
};
\end{axis}
\end{tikzpicture}
\caption{Simulated daily log returns exhibiting volatility clustering.
Calm periods (blue, days 1--100, 221--350, 421--500) have returns within $\pm 0.5\%$;
turbulent periods (red, days 101--220, 351--420) reach $\pm 3\%$.
The clustering is visible by eye: large returns beget large returns.}
\label{fig:vol-clustering-ts}
\end{figure}

Figure~\ref{fig:acf-comparison} quantifies this.
Return autocorrelations hover near zero at all lags, but squared-return autocorrelations remain positive for many lags, confirming that volatility is persistent.

\begin{figure}[htbp]
\centering
\begin{tikzpicture}
\begin{axis}[
  name=acf_returns,
  width=0.48\textwidth,
  height=5.5cm,
  xlabel={Lag (days)},
  ylabel={Autocorrelation},
  xmin=0, xmax=22,
  ymin=-0.15, ymax=0.45,
  grid=major,
  grid style={gray!20},
  xtick={1,5,10,15,20},
  tick label style={font=\small},
  label style={font=\small},
  title style={font=\small\bfseries},
  title={Returns $r_t$},
  ybar,
  bar width=4pt,
]
% Near-zero autocorrelations for returns
\addplot[fill=defblue, draw=defblue!80] coordinates {
(1,0.03)(2,-0.02)(3,0.01)(4,-0.04)(5,0.02)(6,0.01)(7,-0.03)(8,0.02)(9,-0.01)(10,0.03)
(11,-0.02)(12,0.01)(13,-0.01)(14,0.02)(15,-0.03)(16,0.01)(17,0.02)(18,-0.02)(19,0.01)(20,-0.01)
};
% Significance bands
\addplot[dashed, warnred, thick, mark=none] coordinates {(0,0.063)(22,0.063)};
\addplot[dashed, warnred, thick, mark=none] coordinates {(0,-0.063)(22,-0.063)};
\end{axis}

\begin{axis}[
  at={(acf_returns.east)},
  anchor=west,
  xshift=1.2cm,
  width=0.48\textwidth,
  height=5.5cm,
  xlabel={Lag (days)},
  ylabel={Autocorrelation},
  xmin=0, xmax=22,
  ymin=-0.15, ymax=0.45,
  grid=major,
  grid style={gray!20},
  xtick={1,5,10,15,20},
  tick label style={font=\small},
  label style={font=\small},
  title style={font=\small\bfseries},
  title={Squared Returns $r_t^2$},
  ybar,
  bar width=4pt,
]
% Strong, slowly decaying autocorrelations for squared returns
\addplot[fill=keyorange, draw=keyorange!80] coordinates {
(1,0.38)(2,0.31)(3,0.27)(4,0.24)(5,0.21)(6,0.19)(7,0.17)(8,0.15)(9,0.14)(10,0.13)
(11,0.12)(12,0.11)(13,0.10)(14,0.09)(15,0.09)(16,0.08)(17,0.08)(18,0.07)(19,0.07)(20,0.06)
};
% Significance bands
\addplot[dashed, warnred, thick, mark=none] coordinates {(0,0.063)(22,0.063)};
\addplot[dashed, warnred, thick, mark=none] coordinates {(0,-0.063)(22,-0.063)};
\end{axis}
\end{tikzpicture}
\caption{Autocorrelation functions for returns (left) and squared returns (right), based on representative daily equity index data.
Returns show no significant autocorrelation at any lag (all bars within the dashed 95\% significance bands).
Squared returns show strong autocorrelation that decays slowly, reflecting volatility clustering.}
\label{fig:acf-comparison}
\end{figure}

\subsection{Fact 3: Fat Tails}
\label{sec:fact-fat-tails}

If returns were normally distributed, a daily move of 4 standard deviations would occur about once every 126 years.
In practice, moves this large happen several times per year.
The return distribution has ``fat tails'' (also called heavy tails): extreme returns are far more frequent than a Gaussian model predicts.

\citet{Mandelbrot1963} first documented this for cotton prices, proposing that returns follow a stable Paretian distribution rather than a Gaussian.
\citet{Fama1965} confirmed the finding for stock returns, observing leptokurtic (heavy-tailed) distributions across U.S.\ equities.

\begin{definition}[Kurtosis]
The kurtosis of a distribution measures the weight of its tails relative to a normal distribution:
\begin{equation}
  \kappa = \frac{\E\bigl[(r_t - \mu)^4\bigr]}{\bigl(\E\bigl[(r_t - \mu)^2\bigr]\bigr)^2}
  \label{eq:kurtosis}
\end{equation}
\begin{itemize}[nosep]
  \item $\kappa$: kurtosis
  \item $\mu = \E[r_t]$: population mean
  \item For a normal distribution, $\kappa = 3$
  \item \emph{Excess kurtosis} $= \kappa - 3$; values above zero indicate fatter tails than normal
\end{itemize}
Typical daily equity index returns have excess kurtosis in the range 5--10 \citep{Cont2001}.
\end{definition}

Figure~\ref{fig:fat-tails-qq} illustrates fat tails with a Q--Q plot.
If the data were normal, all points would lie on the diagonal.
The S-shaped departure shows that extreme returns occur more often than a normal distribution allows.

\begin{figure}[htbp]
\centering
\begin{tikzpicture}
\begin{axis}[
  width=0.55\textwidth,
  height=7cm,
  xlabel={Normal theoretical quantiles},
  ylabel={Sample quantiles (daily returns, \%)},
  xmin=-4, xmax=4,
  ymin=-6, ymax=6,
  grid=major,
  grid style={gray!20},
  tick label style={font=\small},
  label style={font=\small},
  title style={font=\small\bfseries},
  title={Q--Q Plot: Daily Returns vs.\ Normal Distribution},
]
% Reference line (perfect normal)
\addplot[dashed, gray, thick, mark=none, domain=-4:4] {x*1.1};

% S-shaped Q-Q data showing fat tails
% Lower tail: points fall below the line (more extreme than normal)
\addplot[only marks, mark=*, mark size=1.2pt, defblue] coordinates {
(-3.5,-5.8)(-3.2,-5.1)(-3.0,-4.6)(-2.8,-4.1)(-2.6,-3.7)
(-2.5,-3.4)(-2.3,-3.1)(-2.2,-2.9)(-2.1,-2.7)(-2.0,-2.5)
(-1.9,-2.35)(-1.8,-2.2)(-1.7,-2.05)(-1.6,-1.9)(-1.5,-1.75)
(-1.4,-1.6)(-1.3,-1.48)(-1.2,-1.36)(-1.1,-1.24)(-1.0,-1.12)
(-0.9,-1.01)(-0.8,-0.89)(-0.7,-0.78)(-0.6,-0.67)(-0.5,-0.55)
(-0.4,-0.44)(-0.3,-0.33)(-0.2,-0.22)(-0.1,-0.11)(0.0,0.0)
(0.1,0.11)(0.2,0.22)(0.3,0.33)(0.4,0.44)(0.5,0.55)
(0.6,0.67)(0.7,0.78)(0.8,0.89)(0.9,1.01)(1.0,1.12)
(1.1,1.24)(1.2,1.36)(1.3,1.48)(1.4,1.6)(1.5,1.75)
(1.6,1.9)(1.7,2.05)(1.8,2.2)(1.9,2.35)(2.0,2.5)
(2.1,2.7)(2.2,2.9)(2.3,3.1)(2.5,3.4)(2.6,3.7)
(2.8,4.1)(3.0,4.6)(3.2,5.1)(3.5,5.8)
};

% Annotations
\node[anchor=west, font=\small, warnred] at (axis cs: 2.4, 5.5) {Upper tail: heavier};
\node[anchor=west, font=\small, warnred] at (axis cs: 2.4, 5.0) {than normal};
\draw[-{Stealth}, warnred, thick] (axis cs: 2.3, 5.2) -- (axis cs: 3.3, 5.6);

\node[anchor=east, font=\small, warnred] at (axis cs: -2.3, -5.5) {Lower tail: heavier};
\node[anchor=east, font=\small, warnred] at (axis cs: -2.3, -6.0) {than normal};
\draw[-{Stealth}, warnred, thick] (axis cs: -2.4, -5.3) -- (axis cs: -3.3, -5.6);

\node[anchor=south, font=\small, gray] at (axis cs: 0.8, -1.5) {Normal reference line};
\end{axis}
\end{tikzpicture}
\caption{Q--Q plot of representative daily equity returns against a normal distribution.
If returns were Gaussian, the dots would lie on the dashed diagonal.
The S-shaped departure, first documented by \citet{Fama1965}, shows that both tails of the empirical distribution are heavier than normal: extreme moves (positive and negative) occur much more often than a Gaussian model predicts.}
\label{fig:fat-tails-qq}
\end{figure}

\begin{warning}[Why Fat Tails Matter for Models]
Any model that assumes normally distributed returns will underestimate the probability of large moves.
This affects risk management (VaR breaches), option pricing (mispricing of out-of-the-money options), and backtesting (understating drawdowns).
Chapters~\ref{ch:jumps} and \ref{ch:garch} introduce models that accommodate fat tails.
\end{warning}

\subsection{Fact 4: The Leverage Effect}
\label{sec:fact-leverage}

Negative returns tend to increase future volatility more than positive returns of the same magnitude.
\citet{Black1976} first noted this asymmetry and proposed a mechanism: when a stock price falls, the firm's equity value shrinks relative to its debt, so leverage (debt/equity) rises, making the stock riskier and more volatile.

\begin{definition}[Leverage Effect]
The leverage effect is the negative correlation between an asset's return and changes in its subsequent volatility:
\begin{equation}
  \operatorname{Corr}(r_t,\, \sigma_{t+1}^2) < 0
  \label{eq:leverage-effect}
\end{equation}
\begin{itemize}[nosep]
  \item $r_t$: return at time $t$
  \item $\sigma_{t+1}^2$: conditional variance at time $t+1$
  \item The negative sign means: price drops $\to$ volatility rises
\end{itemize}
\end{definition}

In equity markets, the leverage effect is pronounced.
In foreign exchange and commodity markets, the asymmetry is weaker or sometimes reversed \citep{Cont2001}.
Chapter~\ref{ch:garch} introduces GJR-GARCH and EGARCH models that capture this asymmetry directly.

\subsection{Summary of Stylized Facts}

\begin{keyidea}[Four Facts That Shape Every Volatility Model]
\begin{enumerate}[nosep]
  \item \textbf{Uncorrelated returns:} tomorrow's return direction is unpredictable from today's.
  \item \textbf{Volatility clustering:} large moves follow large moves, small follow small.
  \item \textbf{Fat tails:} extreme returns occur far more often than a Gaussian model predicts.
  \item \textbf{Leverage effect:} negative returns raise future volatility more than positive returns.
\end{enumerate}
Any useful volatility model must reproduce (at minimum) facts 2 and 3.
A good model also captures fact 4.
\end{keyidea}

% ═══════════════════════════════════════════════════
\section{Conditional vs.\ Unconditional Volatility}
\label{sec:cond-uncond}
% ═══════════════════════════════════════════════════

Everything so far has treated volatility as a single number: the standard deviation of the full return sample.
That number is the \emph{unconditional} volatility.
It answers the question ``what is the typical size of a daily return, averaging over all market conditions?''

Traders need a different answer: given what has happened up to today, how volatile will the market be \emph{tomorrow}?
That is the \emph{conditional} volatility.

\begin{definition}[Unconditional vs.\ Conditional Variance]
\textbf{Unconditional variance:}
\begin{equation}
  \sigma^2 = \operatorname{Var}(r_t)
  \label{eq:uncond-var}
\end{equation}
A single number, constant over time. It is the long-run average variance.

\textbf{Conditional variance:}
\begin{equation}
  \sigma_t^2 = \operatorname{Var}(r_t \mid \mathcal{F}_{t-1})
  \label{eq:cond-var}
\end{equation}
\begin{itemize}[nosep]
  \item $\sigma_t^2$: variance of the return at time $t$, conditional on past information
  \item $\mathcal{F}_{t-1}$: the information set (filtration) available at time $t-1$, which includes all past returns, prices, and any other observable data up through time $t-1$
\end{itemize}
The conditional variance changes from day to day as new information arrives.
\end{definition}

\begin{intuition}[Weather Analogy]
Unconditional volatility is like the average annual rainfall for your city: a useful summary, but it does not tell you whether to carry an umbrella tomorrow.
Conditional volatility is the weather forecast: it uses recent data (yesterday's pressure, today's cloud cover) to predict tomorrow's conditions.
This guide is about building the forecast, not computing the average.
\end{intuition}

The unconditional variance is related to the conditional variance by the law of total variance:
\begin{equation}
  \sigma^2 = \E[\sigma_t^2] + \operatorname{Var}\bigl(\E[r_t \mid \mathcal{F}_{t-1}]\bigr)
  \label{eq:total-variance}
\end{equation}
\begin{itemize}[nosep]
  \item $\E[\sigma_t^2]$: average of the conditional variances over time
  \item $\operatorname{Var}\bigl(\E[r_t \mid \mathcal{F}_{t-1}]\bigr)$: variance of the conditional mean (small for equities, since expected returns are nearly constant at daily frequency)
\end{itemize}

When the conditional mean is approximately constant (a reasonable assumption for daily equity returns), the unconditional variance is simply the time average of the conditional variances.
The gap between the two is what makes volatility modeling interesting.

\begin{keyidea}[The Central Goal of This Guide]
The goal of every model in Chapters~\ref{ch:garch}--\ref{ch:hybrid} is to produce accurate estimates of the conditional variance $\sigma_t^2$.
Chapter~\ref{ch:rv} shows how to measure the realized conditional variance from high-frequency data.
Chapters~\ref{ch:garch}--\ref{ch:hybrid} show how to forecast it.
Chapter~\ref{ch:evaluation} shows how to tell whether your forecasts are any good.
\end{keyidea}

% ═══════════════════════════════════════════════════
% SUMMARY
% ═══════════════════════════════════════════════════

\section*{Summary}
\addcontentsline{toc}{section}{Summary}

\begin{itemize}[nosep]
  \item \textbf{Returns} convert raw prices into comparable, analyzable quantities.
  Log returns ($r_t = \ln P_t - \ln P_{t-1}$) are preferred because they are additive over time.

  \item \textbf{Simple vs.\ log returns} are nearly identical for small magnitudes (daily equities); they diverge for large moves.

  \item \textbf{Sample variance} ($\hat{\sigma}^2$) measures dispersion around the mean return; standard deviation ($\hat{\sigma}$) is its square root.

  \item \textbf{Annualized volatility} scales daily volatility by $\sqrt{252}$, assuming independent returns.

  \item \textbf{Volatility matters} because it is a direct input to options pricing, risk management (VaR), portfolio optimization, and trade execution.

  \item \textbf{Stylized fact 1:} Returns are approximately uncorrelated (no predictability in the mean).

  \item \textbf{Stylized fact 2:} Squared and absolute returns are strongly autocorrelated (volatility clusters).

  \item \textbf{Stylized fact 3:} Return distributions have fat tails (excess kurtosis of 5--10 for daily equity returns).

  \item \textbf{Stylized fact 4:} The leverage effect makes volatility respond asymmetrically to negative vs.\ positive returns.

  \item The \textbf{unconditional} variance is a single long-run number; the \textbf{conditional} variance $\sigma_t^2$ varies over time and is what we want to forecast.

  \item The \textbf{central goal} of this guide is to estimate and forecast the conditional variance $\sigma_t^2$ using both classical and ML methods.

  \item \citet{Cont2001} provides the canonical reference for stylized facts across asset classes.

  \item \citet{Mandelbrot1963} and \citet{Fama1965} established that financial returns have fat tails, not Gaussian distributions.

  \item \citet{Engle1982} introduced ARCH, the first formal model of time-varying conditional variance, covered in Chapter~\ref{ch:garch}.
\end{itemize}

% ═══════════════════════════════════════════════════
% KEY RESULTS TABLE
% ═══════════════════════════════════════════════════

\begin{table}[htbp]
\centering
\caption{Key results referenced in this chapter.}
\label{tab:ch01-key-results}
\begin{tabular}{@{} p{4.5cm} p{6cm} p{4cm} @{}}
\toprule
\textbf{Paper} & \textbf{Result} & \textbf{Relevance} \\
\midrule
\citet{Mandelbrot1963} &
Daily commodity returns follow a stable Paretian (fat-tailed) distribution, not Gaussian; sample variance does not converge. &
First evidence that normal-distribution assumptions fail for financial data. \\[6pt]

\citet{Fama1965} &
Confirmed fat tails (leptokurtosis) in daily U.S.\ stock returns via S-shaped Q--Q departures from normality. &
Extended Mandelbrot's finding to equities; supported the stable distribution hypothesis. \\[6pt]

\citet{Black1976} &
Documented negative correlation between stock returns and subsequent volatility changes (the leverage effect). &
Motivates asymmetric volatility models (GJR-GARCH, EGARCH) in Chapter~\ref{ch:garch}. \\[6pt]

\citet{Engle1982} &
Introduced ARCH: conditional variance is a function of past squared residuals; applied to UK inflation. &
Founded the entire field of conditional volatility modeling (Chapter~\ref{ch:garch}). \\[6pt]

\citet{Cont2001} &
Canonical enumeration of stylized facts (absence of return autocorrelation, volatility clustering with slow power-law decay, fat tails, leverage effect) across equities, FX, and commodities. &
Benchmark reference for the empirical properties any volatility model must match. \\
\bottomrule
\end{tabular}
\end{table}
